\clearpage
\section{Convolutional Neural Networks (CNN) for Endoscopic Vision}
\label{sec:cnn}

\subsection*{Introduction}
Convolutional Neural Networks (\textbf{CNNs}) are deep learning models designed to process grid-like data, especially images. Their key idea is to learn local visual patterns through convolutional filters and to reuse the same weights across the input space. This gives CNNs strong performance on tasks that involve spatial structure.

CNNs are one of the most important ideas in modern computer vision, and they are used a lot in image classification, feature extraction, segmentation, and medical imaging \cite{lecun1998cnn,goodfellow2016deep,litjens2017dlSurvey}.

Some of the biggest milestones in CNNs are \textbf{AlexNet}, \textbf{ResNet}, and \textbf{U-Net} \cite{krizhevsky2012alexnet,he2016resnet,ronneberger2015unet}.

\subsection{Core Building Blocks}
The typical CNN pipeline consists of convolution layers, activation functions, pooling layers, and one or more fully connected layers at the end. Each part serves a distinct purpose:

\begin{itemize}
\item \textbf{Convolution:} Extracts local features such as edges, textures, and shapes.
\item \textbf{Activation:} Adds non-linearity, commonly using ReLU.
\item \textbf{Pooling:} Reduces spatial size and improves robustness to small translations.
\item \textbf{Dense Layers:} Combine extracted features to produce final predictions.
\end{itemize}

\begin{table}[h]
\centering
\caption{Main components of a CNN}
\begin{tabular}{|l|l|l|}
\hline
\textbf{Component} & \textbf{Function} & \textbf{Effect} \\ \hline
Convolution layer & Feature extraction & Detects local patterns \\ \hline
ReLU activation & Non-linearity & Improves expressive power \\ \hline
Pooling layer & Downsampling & Reduces feature map size \\ \hline
Flatten layer & Reshaping & Connects to dense layers \\ \hline
Fully connected layer & Classification/regression & Produces final output \\ \hline
\end{tabular}
\label{tab:cnn_components}
\end{table}

\begin{figure}[h]
\centering
\fbox{\parbox[b][3.5cm][c]{0.9\linewidth}{\centering Placeholder: CNN architecture diagram (conv → pool → FC)}}
\caption{CNN architecture (placeholder).}
\label{fig:cnn_arch}
\end{figure}

\subsection{Why CNNs Work Well for Vision Tasks}
CNNs are effective because they exploit two key properties of images: local spatial correlation and translation consistency. Instead of connecting every input pixel to every hidden unit, CNNs focus on local receptive fields and share filters across the image. This reduces the number of parameters and improves generalization.

\begin{table}[h]
\centering
\caption{Typical CNN applications}
\begin{tabular}{|l|l|l|}
\hline
\textbf{Application} & \textbf{Input Type} & \textbf{Output} \\ \hline
Image encoding & Natural images & Visual features \\ \hline
Appearance modelling & Images or video frames & Local descriptors \\ \hline
Segmentation & Medical or scene images & Pixel-wise labels \\ \hline
OCR and document analysis & Scanned pages & Text or layout features \\ \hline
Industrial inspection & Product images & Defect features \\ \hline
\end{tabular}
\label{tab:cnn_applications}
\end{table}

\subsection{CNNs for Navigation and Extraction}
In this work, CNNs are used as visual encoders for two kinds of image input: colon wall
images and target polyp images. The goal is not to build a polyp detector. Instead, the
CNN helps the agent understand the scene, keep the colon in view, and support the next
action during navigation or extraction.

For navigation, CNN features can help the agent recognise folds, bends, lighting changes,
and the current camera view. For extraction, they can help describe the polyp shape,
position, and nearby tissue so the agent can place tools more carefully. In both cases,
the CNN acts as a perception front-end for the decision policy.

The basic idea is simple: image input goes in, compact visual features come out, and the
policy uses those features to choose the next action. That is the bridge between vision
and control in this project.

Evaluation typically uses metrics that capture both feature quality and downstream
usefulness:
\begin{figure}[h]
\centering
\fbox{\parbox[b][3.5cm][c]{0.9\linewidth}{\centering Placeholder: colon wall and polyp image input for the agent}}
\caption{Example colon wall and polyp image input for the agent (placeholder).}
\label{fig:endoscopic_example}
\end{figure}

\begin{table}[h]
\centering
\caption{Common evaluation metrics for endoscopic visual encoders}
\begin{tabular}{|l|l|l|}
\hline
\textbf{Task} & \textbf{Metric} & \textbf{What It Measures} \\ \hline
Visual encoding & Representation quality & How well features summarize the image \\ \hline
Navigation support & Stability / latency & Real-time usefulness for control \\ \hline
Extraction support & Robustness & How well features work under lighting and motion changes \\ \hline
Deployment & Latency / stability & Real-time usability, reduced flicker \\ \hline
\end{tabular}
\label{tab:endoscopy_metrics}
\end{table}

Practical constraints include strong appearance variability, specular highlights,
occlusions, motion blur, and changes in colon wall texture. These factors motivate
robust training practices such as data augmentation and careful validation across
different acquisition conditions.

\subsection{Limitations}
Despite their success, CNNs have limitations:

\begin{itemize}
\item They require large annotated datasets for strong performance.
\item They are less natural for long-range temporal dependencies than sequence models.
\item Very deep networks can be expensive to train and deploy.
\end{itemize}

\subsection{Summary}
CNNs are still a foundational architecture for visual recognition tasks and are
especially useful for endoscopic perception modules. In this project, they turn images
of the colon wall and the target polyp into useful features for the agent. Those features
are then fed into higher-level decision systems that handle navigation and extraction.